\section*{Введение}
\addcontentsline{toc}{section}{Введение}

В этой главе собраны основные понятия, на которых держится вся дальнейшая теория: кольца и поля, характеристика,
кольцо многочленов $F[X]$ над полем, расширения полей и их степени, алгебраические и трансцендентные элементы. Теория
групп здесь почти не нужна; основной инструмент~--- деление многочленов с остатком и однозначность разложения на
неприводимые множители в $F[X]$.

\begin{itemize}[leftmargin=*]
\item \emph{Кольца и поля} (\rf{1.0a}--\rf{1.4a}). Простые и неприводимые элементы, характеристика поля и простое
подполе, эндоморфизм Фробениуса $a\mapsto a^p$ в характеристике $p$.
\item \emph{Кольцо многочленов} (\rf{1.5}--\rf{1.10}). Универсальное свойство, деление с остатком, корни, формула
Тейлора, алгоритм Евклида и соотношение Безу с точными оценками степеней и единственностью коэффициентов (в матричной
записи, континуанты, цепные дроби, матрица Сильвестра, интерполяция Лагранжа), идеалы кольца $F[X]$.
\item \emph{Разложение многочленов на множители} (\rf{1.11}--\rf{1.18g}). Рациональные корни, лемма Гаусса, критерий
Эйзенштейна (над $\Z$ и над любым факториальным кольцом), алгоритм разложения в $\Q[X]$ с оценками корней и
коэффициентов, приведение по модулю простого числа и многочлен $X^4-10X^2+1$, неприводимый над $\Q$, но приводимый по
модулю каждого простого; подъём Гензеля для разложений и для корней, обзор практических методов разложения.
\item \emph{Расширения полей} (\rf{1.19}--\rf{1.32d}). Степень расширения и её мультипликативность, подкольцо и подполе,
порождённые подмножеством, поле $F[X]/(f)$, поле корня (\rf{1.27a}--\rf{1.27c}), алгебраические и трансцендентные
элементы, минимальный многочлен; примеры с квадратными корнями (\rf{1.32a}--\rf{1.32d}).
\item \emph{Трансцендентные числа} (\rf{1.33}--\rf{1.34b}). Счётность множества алгебраических чисел, числа Лиувилля.
\item \emph{Построения циркулем и линейкой} (\rf{1.34c}, \rf{1.35}--\rf{1.42b}). Степень расширения как препятствие к
построению: удвоение куба, трисекция угла, квадратура круга, правильные многоугольники; построения с трисектором и со
вставкой.
\item \emph{Алгебраически замкнутые поля} (\rf{1.43}--\rf{1.47c}). Равносильные определения, алгебраическое
замыкание, алгебраическое замыкание поля в большей области целостности, поле алгебраических чисел; историческое
замечание.
\item \emph{Упражнения} (\rf{1-1}, \dots, \rf{1-8}) с подробными решениями.
\end{itemize}

Все кольца в этой главе (и во всём тексте) имеют единицу, а гомоморфизмы колец переводят единицу в единицу.
